%% sections/slide08-summary.tex — SLIDE 8 — Summary & future work
%% =============================================================================
%% CONTENT BLUEPRINT — Summary and Future Work
%% =============================================================================
%% Time budget: 45 seconds.
%%
%% WHAT TO PUT HERE:
%%
%%   TOP — one headline paragraph (3 sentences max)
%%     Sentence 1: What was studied (rephrase the title in plain words).
%%     Sentence 2: What was done.
%%     Sentence 3: What was found (the headline number with units).
%%
%%   MIDDLE — verdict table mapping each objective to its outcome.
%%     Two-column table:
%%       | Objective (from slide 3)        | Status                      |
%%       | O1: measure n₂ of SF59          | met (slide 5, value + unc)  |
%%       | O2: extract ... vs strain       | met (slide 6 R2)            |
%%       | O3: compare with literature     | partial (8/10 refs agree; 2 |
%%       |                                 |  flagged for re-examination)|
%%       | O4: assess X                    | not met (out of scope, ...)|
%%
%%     Verdict labels: \highlight{met} / \highlight{partial} / \highlight{not met}.
%%     Be honest. The examiner will read the thesis.
%%
%%   BOTTOM — Future work bullets.
%%     • 3–4 bullets. Each bullet is ONE specific future direction.
%%     • Pattern: "To extend this to [X] by [method], which [physical insight]."
%%     • Optional: include 1-line resource estimate (compute, time, equipment)
%%       for one of the future directions — shows you've thought about it.
%%
%% LAYOUT TIP:
%%   • Use a \begin{beamercolorbox}[wd=.96\paperwidth]{block body example}
%%     box for the headline paragraph so it is visually separated from the
%%     verdict table and bullets.
%% =============================================================================

\begin{frame}{Summary \& Future Work}
  \begin{beamercolorbox}[wd=.96\paperwidth, sep=8pt]{block body example}
    \textbf{\color{sustBlueDark}Headline.}
    \medskip

    \itshape In this work we \highlight{[verb: studied / measured /
    computed]} \highlight{[system]} \highlight{[phenomenon]} by
    \highlight{[method]}, and found that \highlight{[target quantity]}\,=\,
    \highlight{[value]}\,$\pm$\,[uncertainty]\,[unit] under
    \highlight{[conditions]} — \highlight{[agreement / deviation]} with
    literature.
  \end{beamercolorbox}

  \vspace{0.1cm}
  \small
  \begin{tabular}{@{}p{0.62\linewidth}@{\,}p{0.33\linewidth}@{}}
    \toprule
    \textbf{Objective} & \textbf{Status (slide reference)} \\
    \midrule
    O1 — Measure $n_2$ of SF59 at 632.8 nm & \highlight{met} (slide 5) \\
    O2 — Extract $n_2$ vs.\ input intensity & \highlight{met} (slide 6 R2) \\
    O3 — Compare with 10 published refs & \highlight{partial} (8/10 agree; 2 flagged) \\
    O4 — Assess role of thermal lens & \highlight{not met} (out of time) \\
    \bottomrule
  \end{tabular}

  \vspace{0.18cm}
  \textbf{\color{sustBlueDark}Future directions}
  \begin{itemize}\setlength\itemsep{1pt}
    \item Extend to $[$second material$]$ at $[$temperature range$]$ to
          isolate $[$effect$]$.
    \item Add BSE on top of GW to resolve $[$open question$]$.
    \item Deploy $[$code$]$ to Zenodo (DOI linked in slide 9) so other
          groups can reproduce.
  \end{itemize}
\end{frame}
